\documentclass[11pt]{article}
\usepackage{amsmath,amssymb,amsthm}
\usepackage{graphicx}
\usepackage[margin=1in]{geometry}
\usepackage{booktabs}
\usepackage{hyperref}

\newtheorem{theorem}{Theorem}
\newtheorem{definition}{Definition}

\title{Learning Mathematics Through the Machine: \\
Rediscovering Closed-Form Expressions from Data via the LSSR Pipeline}
\author{Rainne Eliot Busque \and Mikhail \\
\small Advisor: Dr.\ Jimmie Adriazola, Arizona State University}
\date{June 27, 2026}

\begin{document}

\maketitle

\begin{abstract}
We investigate whether a sparse-recovery pipeline can rediscover known closed-form
mathematical expressions directly from numerical data, without ever being given the
governing equation. We construct a Vandermonde-type candidate library $D$ from sampled
data $b$, solve $Dc = b$, and recover the coefficient vector $c$ via a four-step
procedure we call LSSR (Least Squares $\to$ Sparsify $\to$ Refine): least-squares solve,
hard-threshold small coefficients, and refit on the surviving support. Across three
stages of increasing difficulty --- a finite polynomial, the Maclaurin series of $e^{t}$,
and the bivariate family $e^{at}$ --- we identify and isolate three independent failure
modes: underdetermination, truncation, and collinearity. We show that only collinearity
is actually fixed by sparsification; the other two are numerical issues that more data
or a better threshold cannot repair. Two independent implementations (a matrix-form
solve and a reduced vector-form solve at $a=1$) converge on the same diagonal structure
and the same conclusion: QR/SVD-based solves outperform normal equations by roughly nine
orders of magnitude in conditioning, and a fixed relative threshold of $10\%$ silently
destroys true coefficients smaller than $1/3!$.
\end{abstract}

\section{Introduction}

Most uses of machine learning fit a black-box function to data: given inputs and
outputs, the model interpolates a relationship that may have no algebraic meaning beyond
the fit itself. Our project asks a different question. Suppose we already know that the
underlying relationship has a sparse representation in a fixed library of candidate
terms --- monomials in our case. Can an optimization algorithm, given nothing but
samples of the function, recover the exact coefficients of that representation, and
therefore the exact closed-form expression, with no prior knowledge of which terms are
nonzero?

We call this \emph{discovery} rather than \emph{curve-fitting}: a curve fit draws a line
through points, while discovery reads off which terms exist in the underlying expression
and what their exact values are. If those values match known mathematical constants ---
for example, the Taylor coefficients $1/n!$ --- the result is not just an accurate
predictor but a recovered piece of mathematical truth.

Concretely, we build a candidate (Vandermonde) matrix $D$ whose columns are monomial
terms evaluated at sample points, collect a data vector $b$, and seek a coefficient
vector $c$ solving
\begin{equation}
D c = b.
\label{eq:dcb}
\end{equation}
The columns of $D$ define the hypothesis space; the machine is never told which columns
matter. We solve for $c$ via the LSSR pipeline: an ordinary least-squares solve, followed
by hard thresholding of small coefficients, followed by a refit restricted to the
surviving (non-thresholded) columns. The project proposal deliberately specifies this
naive threshold-based scheme rather than an $\ell_1$-penalized method such as LASSO or
ISTA, so that the failure modes of simple sparsification could be characterized on their
own terms before introducing a more principled penalty.

We organize the investigation into three stages of increasing difficulty:

\begin{enumerate}
  \item \textbf{Stage 1 (polynomial):} recover $x(t) = t + t^2$, a finite-degree
  ground truth, as a sanity check on the pipeline itself.
  \item \textbf{Stage 2 ($e^{t}$):} recover the Maclaurin series of $e^{t}$, an
  infinite series that must be approximated by a finite basis, introducing truncation
  error for the first time.
  \item \textbf{Stage 3 ($e^{at}$):} recover the bivariate family $e^{at}$ via a
  $4\times 4$ grid of monomials $a^m t^n$, whose true coefficient matrix is diagonal
  ($A_{mn} = 1/n!$ when $m=n$, zero otherwise). This stage introduces severe collinearity
  between candidate columns.
\end{enumerate}

Across these three stages we found that apparent ``failures'' of the pipeline were not
one problem but three independent ones --- underdetermination, truncation, and
collinearity --- each with a different cause and a different fix, and that only one of
the three (collinearity) is actually addressed by sparsification at all.

\section{Method}

\subsection{The LSSR pipeline}

Given a candidate matrix $D \in \mathbb{R}^{N \times K}$ (samples $\times$ candidate
terms) and data $b \in \mathbb{R}^{N}$, the pipeline proceeds in four steps:

\begin{enumerate}
  \item \textbf{Generate data.} Samples are produced either by numerically integrating
  the governing ODE with Euler's method, or by direct exact evaluation of the closed
  form (e.g.\ \texttt{np.exp(a*t)}), to separate data-generation error from solver error.
  \item \textbf{Least squares.} Solve Equation~\eqref{eq:dcb} for the full,
  dense coefficient vector $c$, typically via \texttt{np.linalg.lstsq}, which uses an
  SVD-based solve rather than forming the normal equations directly.
  \item \textbf{Sparsify.} Apply a hard threshold: any coefficient with
  $|c_j| < \varepsilon$ is set to exactly zero, where $\varepsilon$ is a fixed fraction
  (here, $10\%$) of the largest coefficient magnitude.
  \item \textbf{Refine.} Re-solve the least-squares problem restricted to the
  surviving (non-zeroed) columns of $D$, removing the bias introduced by the discarded
  candidates' having been forced into the original full solve.
\end{enumerate}

\subsection{Two independent implementations}

To guard against a single coding error masquerading as a scientific conclusion, the two
authors implemented the pipeline independently, without viewing one another's code, and
compared results only afterward.

\begin{itemize}
  \item \textbf{Rainne} built the matrix-form pipeline by hand in guided notebooks: a
  Vandermonde-style design matrix solved via \texttt{np.linalg.lstsq} (SVD), carried
  through Stages 1--3 with threshold-and-refit.
  \item \textbf{Mikhail} (\texttt{learn\_exp\_afixed.py}) fixed $a=1$, collapsing the
  problem to a single-variable, degree-11 Taylor recovery of $e^{t}$, solved both via the
  normal equations and via QR/SVD, with BFGS used for refinement.
\end{itemize}

Despite different code, different parameterizations, and different solvers, both
authors independently converged on the same stable SVD-based solution and the same
diagnosis of the underlying numerical issues, which we treat as the strongest validation
available within the scope of this project.

\section{Results}

\subsection{Stage 1: a finite polynomial as a sanity check}

For $x(t) = t + t^2$, data were generated by Euler-integrating
$\dot{x} = 1 + 2t$. Plain least squares alone --- no sparsification required ---
recovered
\[
c_1 = 0.99, \qquad c_2 = 1.00,
\]
with every other candidate coefficient at machine zero ($\sim 10^{-15}$). Because the
ground truth is an exact, finite-degree polynomial, there is no tail to truncate and no
collinearity among the only two active columns, so this stage validates the pipeline's
correctness rather than stressing it.

\subsection{Stage 2: the Maclaurin series of $e^{t}$}

Recovering $e^{t}$ from an 11-term candidate basis returned coefficients matching $1/n!$
to roughly seven digits of accuracy --- visibly $1, 1, \tfrac{1}{2}, \tfrac{1}{6},
\dots$, the Taylor series, obtained with no prior knowledge of the closed form. Three
control experiments isolated the source of the remaining error:

\begin{enumerate}
  \item Replacing Euler-integrated data with exact evaluation did \emph{not} reduce
  the error, showing the error is structural (a property of approximating an infinite
  series with a finite basis) rather than a data-quality issue.
  \item Raising the polynomial degree (4 $\to$ 7) pushed the distortion into terms we
  were already discarding, improving accuracy on the terms we kept.
  \item Shrinking the sampling window produced roughly a tenfold reduction in error,
  because high powers of $t$ shrink fastest when $t$ is restricted to a small interval.
\end{enumerate}

\subsection{Stage 3: the bivariate family $e^{at}$}

Expanding $e^{at}$ in both $a$ and $t$ produces a $4\times 4$ grid of 16 monomial terms
$a^m t^n$. The true coefficient matrix is diagonal, $A_{mn} = 1/n!$ for $m=n$ and zero
otherwise, so 12 of the 16 candidate terms should vanish exactly.

In the original window and degree, this stage failed outright: spurious (``junk'')
coefficients reached magnitudes up to $-1.24$, larger than a genuine diagonal term
($0.95$). No threshold value could separate signal from junk, and switching to exact
(noise-free) data did not help --- direct evidence that the failure is a property of the
matrix $D$ itself (collinearity among its columns), not of the data.

Shrinking the sampling window to $[0, 0.5]$ and raising the degree to 9 resolved this:
with a loosened threshold, LSSR recovered the full diagonal structure to within
$\sim 10^{-3}$ through the $1/5!$ term, confirming the answer is exactly $e^{at}$.

\subsection{The reduced problem: fixing $a=1$}

Mikhail's independent implementation fixed $a=1$, collapsing the $(M+1)\times(N+1)$
coefficient matrix to an $(N+1)$-dimensional vector problem ($e^{t}$ alone) and pushing
the degree to 11. Comparing solvers on identical data:
\begin{itemize}
  \item Normal equations: error $\approx 40$.
  \item QR/SVD: error $\approx 1.2 \times 10^{-7}$.
\end{itemize}
The data were exact, with no injected measurement noise, so this nine-order-of-magnitude
gap is purely a consequence of solver conditioning: forming the normal equations squares
the condition number of $D$, while QR/SVD solves $Dc=b$ directly and preserves
$\mathrm{cond}(D)$. High-order terms eventually hit the floating-point floor
($1/10! \approx 3\times 10^{-7}$), marking the limit of numerical, rather than
statistical, accuracy.

\begin{table}[h]
\centering
\caption{Recovery error by solver, fixed $a=1$, degree 11, identical exact data.}
\label{tab:solver-comparison}
\begin{tabular}{@{}lc@{}}
\toprule
Solver & Recovery error \\
\midrule
Normal equations ($D^{\top}D\, c = D^{\top}b$) & $\approx 40$ \\
QR / SVD (direct solve of $Dc=b$)              & $\approx 1.2 \times 10^{-7}$ \\
\bottomrule
\end{tabular}
\end{table}

\section{Three Independent Failure Modes}

The results above show that what initially looked like one fragile pipeline is better
understood as three separate, independently-caused failure modes, only one of which
sparsification actually fixes.

\begin{definition}[Underdetermination]
A system $Dc = b$ with more unknown coefficients than independent equations has no
unique solution. In our setting this is resolved simply by collecting enough samples;
empirically, accuracy saturates once the sample count reaches roughly 30--100, after
which additional samples are wasted computation.
\end{definition}

\begin{definition}[Truncation]
When the ground truth is an infinite series and the candidate library is a finite-degree
basis, the discarded tail
\[
e^{t} - P_N(t) = \sum_{n=N+1}^{\infty} \frac{t^n}{n!}
\]
is unavoidable. The error floor scales as $|{\rm tail}| \approx T_{\max}^{N+1}/(N+1)!$,
which links degree $N$ and sampling window $T_{\max}$: they must be tuned together. A
narrow window starves high-order terms of signal; a wide window with low degree lets the
tail explode. This failure mode was invisible in Stage 1, where the ground truth was an
exact finite polynomial with a zero tail.
\end{definition}

\begin{definition}[Collinearity]
When candidate columns of $D$ are nearly linearly dependent, coefficient weight smears
unpredictably across columns that should be zero. Critically, this is a property of the
\emph{design matrix}, not of the data: exact, noise-free data does not repair it. The
only effective remedy is sparsification (LSSR's threshold-and-refit), not more samples.
\end{definition}

\begin{theorem}
For the Stage~3 design matrix, normal-equation and QR/SVD solves of the identical
least-squares problem $Dc=b$ are mathematically equivalent at infinite precision but
differ by roughly nine orders of magnitude in floating-point practice, because forming
$D^{\top}D$ squares the condition number of $D$.
\end{theorem}
\begin{proof}
Let $D$ have singular values $\sigma_1 \ge \cdots \ge \sigma_k > 0$, so that
$\mathrm{cond}(D) = \sigma_1/\sigma_k$. The matrix $D^{\top}D$ has eigenvalues
$\sigma_i^2$, so $\mathrm{cond}(D^{\top}D) = \sigma_1^2/\sigma_k^2 = \mathrm{cond}(D)^2$.
Solving the normal equations $D^{\top}Dc = D^{\top}b$ inherits this squared condition
number, amplifying floating-point error roughly in proportion to it, whereas QR and SVD
factor $D$ directly and never form $D^{\top}D$, so error amplification scales only with
$\mathrm{cond}(D)$ itself. The empirical gap reported in Table~\ref{tab:solver-comparison}
($\sim 10^{-7}$ versus $\approx 40$) is consistent with this squaring relationship at the
condition numbers present in our Stage~3 / fixed-$a$ design matrices.
\end{proof}

\section{Discussion}

A second, related result concerns the threshold step itself. The true coefficients in
the $e^{t}$ recovery span many orders of magnitude: $1, 1, \tfrac12, \dots,
1/11! \approx 2.5\times 10^{-8}$. The proposal's fixed relative rule --- zero any
coefficient below $10\%$ of the largest --- annihilates every true coefficient past
$1/3!$, because a relative cutoff cannot distinguish a small \emph{true} coefficient from
a small \emph{spurious} one. Loosening $\varepsilon$ to preserve the small true terms
recovers all twelve correctly, but at the cost of sparsifying nothing at all: the
pipeline collapses back to plain least squares plus a refit. This is, in our view, the
sharpest practical limitation of naive hard thresholding relative to a principled
$\ell_1$ penalty (LASSO/ISTA), which shrinks all coefficients continuously rather than
applying one global cutoff.

Finally, the comparison between the full bivariate matrix form and Mikhail's reduced
vector form ($a=1$) shows that dimensionality reduction, where it is legitimate, is worth
far more than additional data or a better solver. The vector form wins on three
independent counts: fewer unknowns, no all-zero off-diagonal block available to absorb
numerical error, and a design matrix that is roughly $10^{8}\times$ better conditioned.
The caveat is that the full matrix form solves a strictly harder and more general
problem --- it must discover, rather than assume, that $a$ and $t$ enter the true
function only through the product $at$.

\section{Conclusion}

Sparse recovery can rediscover known closed-form mathematics directly from sampled data,
with no governing equation supplied in advance: across three stages we recovered an
exact polynomial, the Taylor series of $e^{t}$, and the full diagonal structure of
$e^{at}$, the last validated by two independent implementations built without
cross-communication. The central finding is that what appears to be a single fragile
pipeline is in fact three independent failure modes --- underdetermination, truncation,
and collinearity --- each diagnosed by a different controlled experiment and fixed by a
different remedy (more samples; coupling degree and window; sparsification,
respectively). Of the three, only collinearity is addressed by sparsification at all,
and the naive fixed-threshold version of sparsification used here is itself fragile: it
silently destroys true small coefficients unless loosened to the point of doing nothing.
This motivates replacing the hard threshold with a principled $\ell_1$ penalty solved via
ISTA, a direction for the next phase of this work.

\end{document}
